\subsection{残る三型と再開する具体的な入口}
\label{ns25:support-agenda}\label{ns24:support-agenda}\label{ns23:support-agenda}\label{ns22:support-agenda}\label{ns20:support-agenda}

再開時は最初に \nolinkurl{00_CONTINUATION_POLICY_ja.txt} を読む．
現在の高分母支持の必要候補は\emph{三型・二十六標識}である．
全候補は五頂点を使い，三点連結，最小次数三以上，八高辺以上，$b_1\ge4$ となる．
単位枝は高々二辺の matching であり，同じ core に二本が接すれば全分母一となる．
補 graph の孤立頂点を省略すると，候補は次の通りである．
\begin{center}\small
\begin{tabular}{c|c|c|c|r}
型&高辺数&次数列&低辺補 graph&標識数\\\hline
$G_{5,14}$&8&$(43333)$&$2K_2$&15\\
$G_{5,16}$&9&$(44433)$&$K_2$&10\\
$G_{5,17}$&10&$(44444)$&$\varnothing$&1
\end{tabular}
\end{center}
識別子は詰め直さない．検算は \nolinkurl{verify_support_agenda_v23.py} である．
候補数は整数接着，atlas，幾何的非同型類の存在数ではない．

\paragraph{高分母全体の原始性．}
非空 $H$ には八辺以上ある．全高分母を割る素数があるとすると，素数支持が $H$ と一致するため，
密な素数支持の定理はその素数を五，支持を $K_{2,2,1}$ に強制する．
補 graph の二辺は実際には単位分母であるが，合同分類はその和を零 modulo five にするため矛盾する．
従って高分母全体の最大公約数は一である．

二本の単位枝がある場合には任意七本の分母の最大公約数も一となり，
$G_{5,14}$ で五の倍数の分母は高々六本である．
密な $G_5=K_{2,2,1}$ と一本の実単位枝が共存する場合は $H=G_{5,16}$ に限られ，
五進支持の外のもう一本は単位枝と頂点を共有せず，分母 $q\equiv9\pmod{10}$，$q\ge9$ となる．
合同条件だけで整数 lifting を保証してはいない．

\paragraph{優先候補：二つの四node単位枝と中央 matching．}
Book，$G_{5,15}$，四高 $P_4$，$G_{5,14}$ の五node sector は解決済みであり再探索しない．
$G_{5,14}$ の単位枝を $01,23$ とすれば，両方の配分は $\{4,1,0\}$，
core $4$ の正 matching は $\{01,23\}$ に限られる．
各単位枝の共通負位置には分母三が隣接するが，二つの証人が同一辺となる可能性を残す．
中心の二つの正 pair と四本の高枝の node・対向 $D$ を同時に用い，
自由尺度を持つ node 四の正族と反対符号の小例外を全て保って消去する．
$G_{5,16}$ の五node sector には第二単位枝がないので，今回の五node排除を移さない．
固定小分母を含む三高 core の全族分類，または四高 core との元の行和・pair 消去を候補とする．

\paragraph{三高係数の残余．}
$\delta=\pm1$ と比例領域の $x>0$ は全排除済みである．
零軸 $x=y=0,z\ge1,s<t,0\le n\le z$ の端点では
\[
 n=0:\ s\ge3,\ s<b<3s/2,\ z<s/2;\qquad
 n=z:\ s\ge2,\ z\le s,\ 1\le C_z/(b+1)\le s-1.
\]
固定 $s$ で有限になった端点の一般整数根問題と，内部 $0<n<z$ を分ける．
正整数商 $e=C_n/(b+1)$ の四次式を使い，内部の符号が変わる係数を先に調べる．
canonical な向きを保ち，標識交換で $s<t$ を消してはならない．

退化面 $a=0$ では $M(t-y)=tT$，$hM=(b-1)L$ という正整数積を使う．
$x\ge1$ なら $h\ge b$，$1\le T\le M\le b-2$，$\gcd(h,b-1)\ge2$ である．
固定 $b$ には $b+1\le s\le b^2-1$ の有限復元がある．
$x=0$ も含む全退化面には固定 $s$ の既存手順を残す．
一般の $|\delta|\ge2$ では，尺度正値性の格子幅と元残余を同時に扱う．
各固定入力の有限性を入力自体への一様上界や全域の不存在に読み替えない．

\paragraph{幾何への接続と方向転換．}
全ての非零四ray core の正符号を三つの標識付き matching に絞り，
各辺から有理 weights を構成する．別々の core の weights が大域に一致することは追加問題である．
整数候補が残れば全 core の basin areas，integer framings，作用切片を確認する．
五次・六次 atlas の同値，不変量，別の実際の退化，open-corridor surgery への転換も，
研究全体の面白さとインパクトに応じて判断する．
